\chapter{Tensors as Multilinear Maps}

Linearity expresses that a map respects the algebraic structure of a vector space.
Many constructions in mathematics, however, depend on functions that take several
vector arguments simultaneously while remaining linear in each argument when the
others are held fixed. This leads naturally to the notion of multilinearity.

\medskip

Let \(E_1, E_2, \dots, E_n\) and \(F\) be vector spaces over the same field \(K\).

\begin{definition}[Multilinear map]
A map
\[
T : E_1 \times E_2 \times \cdots \times E_n \to F
\]
is called \emph{multilinear} if it is linear in each argument separately.
That is, for each index \(i \in \{1,\dots,n\}\), and for all vectors
\[
x_i, y_i \in E_i, \quad \alpha \in K,
\]
we have
\[
T(x_1,\dots,x_i +_{E_i} y_i,\dots,x_n)
=
T(x_1,\dots,x_i,\dots,x_n)
+_F
T(x_1,\dots,y_i,\dots,x_n),
\]
and
\[
T(x_1,\dots,\alpha \circ x_i,\dots,x_n)
=
\alpha \circ T(x_1,\dots,x_i,\dots,x_n),
\]
with all other arguments held fixed.
\end{definition}

When \(n=1\), a multilinear map is simply a linear map.
Thus multilinear maps are not a new kind of object, but a direct generalization
of linear transformations.

\medskip

Multilinear maps arise naturally in many contexts:
bilinear products, pairings between vectors and linear functionals,
and higher-order interactions that cannot be reduced to a single input.
Their systematic study leads to the concept of tensors.

\medskip

Informally, a tensor is an object designed to encode multilinear behavior
in an intrinsic and coordinate-free way.

\medskip

There are two equivalent viewpoints that will be used throughout this document.
The first treats tensors directly as multilinear maps.

\medskip

Let \(E\) be a vector space over \(K\), and let \(E^*\) denote its dual space,
that is, the vector space of linear maps from \(E\) to \(K\).

\begin{definition}[Tensor as a multilinear map]
A \emph{tensor of type \((k,\ell)\)} on \(E\) is a multilinear map
\[
T : \underbrace{E^* \times \cdots \times E^*}_{k\ \text{times}}
\times
\underbrace{E \times \cdots \times E}_{\ell\ \text{times}}
\;\longrightarrow\; K.
\]
\end{definition}

Such a tensor takes \(k\) linear functionals and \(\ell\) vectors as inputs,
and produces a scalar.
Linearity is required in each argument separately.

\medskip

This definition makes precise several familiar cases:
\begin{itemize}
\item A scalar in \(K\) can be viewed as a tensor of type \((0,0)\).
\item A vector in \(E\) corresponds to a tensor of type \((0,1)\).
\item A linear functional in \(E^*\) is a tensor of type \((1,0)\).
\item A bilinear form on \(E\) is a tensor of type \((0,2)\).
\item If $E$ is finite-dimensional, a linear map $T : E \to E$ can be identified
      with a tensor of type $(1,1)$.
\end{itemize}

Thus vectors and matrices reappear as the lowest-order cases of a single
unified framework.

\medskip

At this level of abstraction, no coordinates, arrays, or geometric interpretations
are involved.
A tensor is defined entirely by how it acts on vectors and linear functionals,
and by the multilinearity of this action.

% CORRECTION 1: The motivation for the tensor product is now clearly separated
% from its definition. The previous version mixed motivation and solution in the
% same block, with the final sentence anticipating the definition about to follow
% (circular). The motivation is now self-contained: it identifies the limitation
% of multilinear maps (no natural composition) and poses the question the tensor
% product will answer, without naming the solution prematurely.

\medskip

The definition above is conceptually clear, but it has a practical limitation.
Linear maps compose: if $S : E \to F$ and $T : F \to G$ are linear, then
$T \circ S : E \to G$ is again linear, and its matrix is the product of the
matrices of $T$ and $S$.
Multilinear maps do not admit an analogous operation.
Their outputs and inputs do not naturally match in a way that allows canonical
composition, and there is no general rule for combining two multilinear maps
into a third while preserving multilinearity.

\medskip

Multilinear maps do form vector spaces under pointwise operations, so they can
be added and scaled.
What is missing is a way to \emph{reduce} multilinearity to linearity: to
replace a multilinear problem by a linear one on a suitably constructed space.
This is precisely what the tensor product achieves.

\medskip

Let \(E_1,\dots,E_n\) be vector spaces over the same field \(K\).
Consider the set of all multilinear maps
\[
T : E_1 \times \cdots \times E_n \longrightarrow F
\]
into an arbitrary vector space \(F\).
The tensor product is defined so that every such multilinear map corresponds
uniquely to a linear map defined on a single vector space.

\medskip

\begin{definition}[Tensor product]
The \emph{tensor product} of the vector spaces \(E_1,\dots,E_n\) is a vector
space, denoted
\[
E_1 \otimes \cdots \otimes E_n,
\]
together with a multilinear map
\[
\otimes : E_1 \times \cdots \times E_n \longrightarrow
E_1 \otimes \cdots \otimes E_n,
\]
satisfying the following universal property:

For every vector space \(F\) and every multilinear map
\[
T : E_1 \times \cdots \times E_n \longrightarrow F,
\]
there exists a unique linear map
\[
\widetilde{T} : E_1 \otimes \cdots \otimes E_n \longrightarrow F
\]
such that
\[
T(x_1,\dots,x_n) = \widetilde{T}(x_1 \otimes \cdots \otimes x_n)
\]
for all \(x_i \in E_i\).
\end{definition}

\medskip

This property characterizes the tensor product completely.
The elements \(x_1 \otimes \cdots \otimes x_n\) are called \emph{simple tensors}.
Every element of the tensor product space is a finite linear combination of
simple tensors.

\medskip

As a consequence, the study of multilinear maps is equivalent to the study of
linear maps defined on tensor products.
Specifically,
\[
\text{Multilinear maps } E_1 \times \cdots \times E_n \to F
\quad\longleftrightarrow\quad
\text{Linear maps } E_1 \otimes \cdots \otimes E_n \to F.
\]

\medskip

This equivalence is the conceptual foundation of tensor theory.
It explains why tensors can be manipulated algebraically, composed with linear
maps, and represented concretely once bases are chosen.

\medskip

Returning to the case of a single vector space \(E\), a tensor of type
\((k,\ell)\) may equivalently be viewed as a linear map
\[
T :
\underbrace{E^* \otimes \cdots \otimes E^*}_{k}
\otimes
\underbrace{E \otimes \cdots \otimes E}_{\ell}
\;\longrightarrow\; K.
\]

\medskip

Thus tensors are linear maps defined on tensor products, and their apparent
complexity arises only when coordinates are introduced.

\medskip

In contrast with the direct sum (Chapter~I), which combines independent
components additively, the tensor product encodes interactions between
components through multilinearity.

\bigskip
\hrule
\medskip

\textit{Note:} Once bases are chosen, tensor products become finite-dimensional
vector spaces, and tensors admit coordinate representations as multidimensional
arrays. These representations depend on the chosen bases and do not define the
tensor itself.

\medskip
\hrule
\bigskip

Let $E$ be a finite-dimensional vector space over $K$, and let
\[
B = \{ e_1, \dots, e_n \}
\]
be a basis of $E$. Its dual basis
\[
B^* = \{ e^1, \dots, e^n \}
\subset E^*
\]
is defined by the relations
\[
e^i(e_j) = \delta^i_j,
\]
where $\delta^i_j$ is the Kronecker delta.

\medskip

Every vector $x \in E$ admits a unique expansion
\[
x = x^i \circ e_i,
\]
and every linear functional $\varphi \in E^*$ admits a unique expansion
\[
\varphi = \varphi_i \circ e^i.
\]

The scalars $x^i$ and $\varphi_i$ are the \emph{components} of $x$ and
$\varphi$ relative to the chosen bases.

\medskip

Let now $T$ be a tensor of type $(k,\ell)$ on $E$, that is,
\[
T : \underbrace{E^* \times \cdots \times E^*}_{k}
\times
\underbrace{E \times \cdots \times E}_{\ell}
\longrightarrow K.
\]

The action of $T$ is completely determined by its values on basis elements:
\[
T(e^{i_1}, \dots, e^{i_k}, e_{j_1}, \dots, e_{j_\ell})
= T^{i_1 \dots i_k}{}_{j_1 \dots j_\ell}.
\]

These scalars are called the \emph{components of $T$ relative to the basis $B$}.
They encode the tensor entirely once a basis has been fixed.

\medskip

Given arbitrary arguments
\[
(\varphi^1, \dots, \varphi^k, x_1, \dots, x_\ell),
\]
multilinearity implies the expansion
\[
T(\varphi^1, \dots, \varphi^k, x_1, \dots, x_\ell)
=
T^{i_1 \dots i_k}{}_{j_1 \dots j_\ell}
\,
\varphi^1_{i_1} \cdots \varphi^k_{i_k}
\,
x_1^{j_1} \cdots x_\ell^{j_\ell}.
\]

Thus index notation is the coordinate expression of multilinearity after a
basis is chosen.

\medskip

Bilinear maps provide the simplest computational example of multilinearity.
A bilinear map
\[
B : V \times W \to U
\]
is linear in each argument separately. Once bases are chosen in the vector
spaces involved, the map is represented by coefficients $t_{ij}^{\;k}$ such
that the components of the output vector satisfy
\begin{equation}
z^{k} = \sum_{i=1}^{m}\sum_{j=1}^{n} t_{ij}^{\;k}\, x^{i} y^{j}.
\end{equation}

The collection of coefficients $(t_{ij}^{\;k})$ forms a rank-three tensor.
This representation shows that familiar computational procedures are coordinate
expressions of multilinear mappings. For example, the multiplication of complex
numbers
\begin{equation}
(x_1 + i x_2)(y_1 + i y_2)
= (x_1 y_1 - x_2 y_2)
+ i(x_1 y_2 + x_2 y_1)
\end{equation}
can be interpreted as the evaluation of two bilinear forms determined by a
$2 \times 2 \times 2$ tensor of coefficients.

From this viewpoint, algorithms manipulating numerical arrays are
implementations of tensorial transformations written in a chosen reference
system. Computation therefore makes explicit the invariant multilinear
structure underlying algebraic operations.

\medskip

In particular, a bilinear map
\[
B : E \times F \to G
\]
between finite-dimensional vector spaces is represented, once bases are chosen,
by components
\[
B^k{}_{ij},
\]
so that for vectors $x \in E$ and $y \in F$,
\[
(B(x,y))^k = B^k{}_{ij}\, x^i y^j.
\]
Thus multidimensional arrays of coefficients appearing in numerical algorithms
are coordinate representations of tensors encoding multilinear operations.

\medskip

Once a basis has been fixed, the components
\[
T^{i_1 \dots i_k}{}_{j_1 \dots j_\ell}
\]
can be arranged in a multidimensional array of scalars.

\medskip

Changing the basis of $E$ changes the numerical values of the components
according to precise transformation rules, while the tensor itself remains
unchanged.

\bigskip
\hrule
\medskip

\textit{Note: Tensors via Transformation Laws}

The definition of a tensor as a multilinear map is intrinsic and
coordinate-free. An equivalent characterization emphasizes how the components
of a tensor transform under a change of basis.

\medskip

This viewpoint makes explicit the invariant content that survives changes of
coordinates, which is precisely what allows tensorial laws to appear in physics
and in engineering applications independent of the chosen representation.

\medskip

Let $B = \{e_i\}$ and $B' = \{e'_i\}$ be two bases of $E$, related by
\[
e'_i = A^j{}_i \, e_j,
\]
where $A^j{}_i$ is an invertible change-of-basis matrix.
The corresponding dual bases satisfy
\[
e'^i = (A^{-1})^i{}_j \, e^j.
\]

\medskip

Let $T$ be a tensor of type $(k,\ell)$ on $E$, with components relative to $B$
given by
\[
T^{i_1 \dots i_k}{}_{j_1 \dots j_\ell}.
\]

The matrix $A^j{}_i$ expresses the new basis vectors in terms of the old ones,
as discussed in Chapter~I.
Accordingly, the transformation rule below describes how the coordinate
representation of $T$ changes when passing from the basis $B$ to $B'$.

\medskip

The components relative to the new basis $B'$ are given by
\[
T'^{i_1 \dots i_k}{}_{j_1 \dots j_\ell}
=
A^{i_1}{}_{p_1}
\cdots
A^{i_k}{}_{p_k}
\,
(A^{-1})^{q_1}{}_{j_1}
\cdots
(A^{-1})^{q_\ell}{}_{j_\ell}
\,
T^{p_1 \dots p_k}{}_{q_1 \dots q_\ell}.
\]

\medskip

Contravariant indices (upper indices) transform with the matrix $A$, while
covariant indices (lower indices) transform with its inverse.
This transformation rule is not an additional assumption.
It is a direct consequence of multilinearity and of how vectors and dual vectors
transform under a change of basis.

\medskip

Conversely, suppose that for each choice of basis one assigns a multidimensional
array of scalars
\[
T^{i_1 \dots i_k}{}_{j_1 \dots j_\ell}
\]
in such a way that these arrays are related by the transformation law above
whenever the basis is changed.
Then these coordinate collections define a unique multilinear map
\[
T :
\underbrace{E^* \times \cdots \times E^*}_{k}
\times
\underbrace{E \times \cdots \times E}_{\ell}
\longrightarrow K,
\]
independent of the chosen basis.

\medskip

Thus the following two viewpoints are equivalent:
\begin{itemize}
\item A tensor is a multilinear map.
\item A tensor is a collection of components that transform according to the
      tensor transformation law.
\end{itemize}

The first definition emphasizes intrinsic structure.
The second emphasizes invariance under change of reference frame.
Their equivalence explains why tensor theory plays a central role in geometry
and physics: tensorial equations remain valid in every coordinate system,
because both sides transform according to the same law.

\medskip
\hrule
\bigskip

Thus:
\begin{itemize}
\item tensors are abstract multilinear objects;
\item arrays are basis-dependent representations;
\item indices record how multilinearity decomposes into coordinates.
\end{itemize}

This distinction is essential for interpreting computational implementations.

\medskip

Up to this point, tensors have been considered as static multilinear objects.
We now introduce the operation that allows them to interact algebraically.

% CORRECTION 3: The broken sentence ("operation. we now introduce...") and the
% repetition of the no-composition remark have both been removed. The transition
% is now a single clean sentence that opens the contraction section.

\medskip

To combine tensors and produce new multilinear objects, one introduces the
fundamental operation of tensor contraction, which generalizes matrix
multiplication and explains the appearance of summation over repeated indices.

\medskip

We denote by
\[
\mathcal{T}^{(k,\ell)}(E)
\]
the vector space of all tensors of type $(k,\ell)$ on $E$.

\medskip

Let $E$ be a finite-dimensional vector space over $K$.
Consider a tensor
\[
T \in \mathcal{T}^{(k,\ell)}(E)
\]
and a tensor
\[
S \in \mathcal{T}^{(m,n)}(E).
\]

Assume that one covariant slot of $T$ and one contravariant slot of $S$ are
selected.
Tensor contraction is the operation that pairs these two slots using the
natural evaluation map
\[
E^* \times E \to K.
\]

\begin{definition}[Tensor contraction]
Let
\[
T \in \mathcal{T}^{(k,\ell)}(E),
\qquad
S \in \mathcal{T}^{(m,n)}(E).
\]
Choosing one covariant slot of $T$ and one contravariant slot of $S$, their
contraction is defined by composing these arguments with the natural evaluation
pairing
\[
\langle\cdot,\cdot\rangle : E^* \times E \to K.
\]
In coordinates, this operation appears as summation over a basis.

\medskip

The resulting tensor has type
\[
(k+m-1,\;\ell+n-1).
\]
\end{definition}

\medskip

Contraction reduces the total number of tensor slots by two: one contravariant
and one covariant slot are paired and eliminated.
No coordinates are required to define this operation.

\medskip

To express contraction in coordinates, a basis must be chosen.
The abstract pairing then appears as summation over repeated indices.

\medskip

Fix a basis $\{e_i\}$ of $E$ and its dual basis $\{e^i\}$.

Let
\[
T^{i}{}_{j}
\quad\text{and}\quad
S^{j}{}_{k}
\]
be the components of two tensors of type $(1,1)$.

\medskip

The contraction over the index $j$ produces a new tensor with components
\[
R^{i}{}_{k}
=
T^{i}{}_{j} \, S^{j}{}_{k}.
\]

Here the repeated index $j$ indicates summation:
\[
R^{i}{}_{k}
=
\sum_{j=1}^{n}
T^{i}{}_{j} \, S^{j}{}_{k}.
\]
This summation is not an additional rule.
It is the coordinate expression of the abstract pairing between $E^*$ and $E$.

\medskip

Let $A : E \to E$ and $B : E \to E$ be linear maps.
Each can be identified with a tensor of type $(1,1)$.
Relative to a basis, their components are written as
\[
A^{i}{}_{j},
\qquad
B^{j}{}_{k}.
\]

The composition $A \circ B : E \to E$ has components
\[
(A \circ B)^{i}{}_{k}
=
A^{i}{}_{j} \, B^{j}{}_{k}.
\]

Thus matrix multiplication is exactly tensor contraction between two $(1,1)$
tensors, and the familiar summation rule for matrix products is a special case
of the abstract pairing between $E^*$ and $E$.

\medskip

At the abstract level we have:
\[
\text{Tensor contraction}
\]
which, after choosing bases, becomes
\[
\text{Index summation}
\quad\longrightarrow\quad
C^{i}{}_{k} = \sum_j A^{i}{}_{j} B^{j}{}_{k}.
\]

% CORRECTION 4: The paragraph beginning "It is fundamental to recognize..."
% has been removed. Its tone was markedly more rhetorical than the surrounding
% text, it introduced the non-standard term "index demolition", and its content
% (contraction as the engine of computation) is already conveyed more precisely
% by the Systolic Arrays note below.

% CORRECTION 2: The discrete grid example (force array) has been moved to here,
% after contraction is defined. The reader now has the algebraic tools to
% understand what the indices of the array mean and how they could be contracted.
% Previously it appeared before contraction, leaving the array without any
% operational context.

\medskip

In computational and modeling contexts, tensors frequently appear through
multidimensional arrays whose indices label independent parameters of a problem.
The following example illustrates how such arrays arise in practice while still
representing tensorial data.

\medskip

Consider a physical system modeled on a discrete two-dimensional grid.
At each spatial point $(i,j)$ and at each time step $t$, a force vector
\[
F(t,i,j) \in \mathbb{R}^2
\]
is measured, with components $(F_x, F_y)$ relative to a fixed spatial frame.

\medskip

Let:
\begin{itemize}
\item $t \in \{1,\dots,T\}$ index time,
\item $(i,j) \in \{1,\dots,N\} \times \{1,\dots,M\}$ index spatial position,
\item $\alpha \in \{1,2\}$ index force components.
\end{itemize}

The collected data can be written as an array
\[
F_{t i j \alpha},
\]
with shape $(T,N,M,2)$.

\medskip

Each index corresponds to an independent modeling choice: time, space, and
vector components.
The array stores numerical values relative to chosen coordinate systems.
Only the component index $\alpha$ corresponds to a genuine vector slot in the
tensorial sense; the indices $(t,i,j)$ label independent parameters of the
model.
Contracting the index $\alpha$ against a linear functional produces, at each
spacetime point, a scalar measurement of the force along a chosen direction.

\medskip

Abstractly, this data can be interpreted as the coordinate representation of a
tensor-valued function on a discrete domain.

\bigskip
\hrule
\medskip

\textit{Note: Linear Algebra Realized in Hardware --- Systolic Arrays
\cite{kung}}

\medskip

The interpretation of matrix multiplication as tensor contraction is not only
a conceptual reformulation.
It also determines how modern computational hardware is designed.

\medskip

In large numerical computations, the dominant operation is the repeated
evaluation of inner products,
\[
c_i = \sum_{j=1}^{n} a_{ij} b_j,
\]
which is the coordinate expression of a linear map applied to a vector.
Equivalently, it is a contraction between a tensor of type $(1,1)$ and a
tensor of type $(0,1)$.

\medskip

In the late 1970s, H.~T.~Kung and C.~E.~Leiserson proposed organizing
computation around this operation through architectures called
\emph{systolic arrays}.
Instead of moving data back and forth between memory and a central processor,
data flows rhythmically through a grid of simple processing elements.
Each element repeatedly performs a single operation of the form
\[
C \leftarrow A B + C,
\]
a multiply-accumulate step implementing one contribution to an inner product.

\medskip

Mathematically, each processing element evaluates a partial tensor contraction.
As data propagates across the array, successive contractions accumulate, and
the global result emerges from the coordinated action of many local linear
operations.

\medskip

This organization reflects a structural principle: hardware efficiency is
achieved when computation mirrors algebraic structure.
Since linear transformations decompose into sums of products, an architecture
optimized for multiply-accumulate operations naturally implements linear algebra
at scale.

\medskip

Modern high-performance computing architectures frequently adopt variants of
this paradigm, arranging many simple processing elements to evaluate
multiply-accumulate operations in parallel.
What appears at the software level as tensor contraction or matrix
multiplication is therefore executed physically as streams of coordinated inner
products.

\medskip

From the present viewpoint, this development is not accidental.
Tensor contraction is the fundamental interaction between multilinear objects,
and large-scale computation reduces complex algorithms to repeated evaluations
of this interaction.

\medskip

In this sense, numerical arrays stored in memory, index summations written in
formulas, and data flowing through silicon circuits are three realizations of
the same mathematical object: the evaluation of multilinear maps.

\medskip
\hrule
\bigskip

The abstract principles developed above appear concretely in classical physics,
where tensorial quantities describe invariant relations between physical
observables.

\medskip

\textit{A classical example: the inertia tensor \cite{feynman}}

In classical mechanics, the distribution of mass in a rigid body is encoded by
a second-order tensor called the \emph{inertia tensor}.
Relative to a chosen orthonormal basis of $\mathbb{R}^3$, its components are
given by
\[
I_{ij} = \int_{\Omega}
\big( \delta_{ij} \|x\|^2 - x_i x_j \big)\, dm,
\]
where $\Omega$ is the body and $x$ is the position vector.

\medskip

The inertia tensor defines a linear relation between angular velocity $\omega$
and angular momentum $L$:
\[
L_i = I_{ij} \, \omega^j.
\]

\medskip

Under a change of basis determined by an invertible matrix $A$, its components
transform according to
\[
I'_{ij}
=
A^{p}{}_{i} \,
A^{q}{}_{j} \,
I_{pq},
\]
which is precisely the transformation law of a tensor of type $(0,2)$.

\medskip

Thus, although the numerical entries of $I_{ij}$ depend on the chosen reference
frame, the geometric object it represents does not.
The physical content---resistance to rotational acceleration---is invariant
under change of coordinates.

\medskip

Tensor theory therefore extends the central principle introduced in Chapter~I:
mathematical structure is identified not by its numerical representation,
but by the invariance of its defining relations under admissible
transformations.

% CORRECTION 5 & 6: The Notes are now reordered to group applied examples
% together and separate them from the speculative outlook.
% Order: Outlook (future hardware) -> chi-square remark (bridge back to Ch. II)
% -> OLAP cubes (data analysis application) -> Systolic arrays (hardware).
% The redundant introductory paragraph that preceded the Systolic Arrays note
% in the main text has been removed; the note's own opening now suffices.

\section*{Notes and Remarks}

\bigskip
\hrule
\medskip

\textit{Outlook: Computation Beyond Electronics \cite{Xu}}

\medskip

Recent developments in computing architectures explore physical substrates
beyond traditional electronic circuits, including optical and photonic systems
in which information is processed through interference, diffraction, and
propagation of light.

\medskip

Despite their physical differences, such systems are designed to implement the
same fundamental algebraic operations encountered throughout this chapter.
In particular, the dominant computational task remains the evaluation of
expressions of the form
\[
C^{i}{}_{k} = A^{i}{}_{j} B^{j}{}_{k},
\]
that is, tensor contraction.

\medskip

In these emerging architectures, the interaction of physical signals naturally
realizes multiply-accumulate operations in parallel, so that the evaluation of
multilinear maps is performed directly by the underlying physical process rather
than by sequential electronic instructions.

\medskip

A non-trivial constraint, however, arises from the physical encoding of
information.
Non-coherent optical systems represent values as light intensities, which are
intrinsically non-negative; they therefore operate within the restricted
domain~$\mathbb{R}^+$ rather than the full real line.
Since tensor contractions require weights and activations that range over all
of~$\mathbb{R}$, such systems are algebraically incomplete for general
computation.
Coherent architectures address this limitation by encoding values as optical
field \emph{amplitudes}---quantities that carry sign---and recovering the
signed output via interference at detection.
The full group $(\mathbb{R},+)$ is thereby made available, and tensor
contractions of arbitrary sign are realizable in hardware.

\medskip

From this perspective, advances in hardware do not alter the mathematical
structure of computation.
They provide new physical realizations of the same invariant operations: linear
transformations and tensor contractions, now extended to their natural domain.
The algebraic completeness required here---the passage from~$\mathbb{R}^+$
to~$\mathbb{R}$---is precisely the structural distinction between a
multiplicative and an additive group, whose correspondence is the subject of
Chapter~IV.

\medskip

Thus, independently of whether computation is carried out by electronic,
optical, or other physical means, the algebraic framework developed here
remains unchanged.

\medskip
\hrule
\bigskip

\bigskip
\hrule
\medskip

\textit{Remark: Bridge to Chapter~II --- the chi-square statistic as a
$(0,2)$ tensor}

\medskip

The chi-square statistic encountered in Chapter~II provides an earlier example
of a tensor of type $(0,2)$.
Recall that the expression
\[
V
=
\sum_{i=1}^{n}
\frac{(Y_i - np_i)^2}{np_i}
\]
was identified as the squared norm induced by a symmetric positive-definite
bilinear form $B$ on $\mathbb{R}^n$, whose matrix in the standard basis is
diagonal with entries $(np_i)^{-1}$.
From the present point of view, $B$ is precisely a tensor of type $(0,2)$: a
multilinear map
\[
B : \mathbb{R}^n \times \mathbb{R}^n \longrightarrow \mathbb{R}
\]
that is symmetric and positive-definite.
Its coordinate representation changes with the basis, but the geometric
object it defines---the ellipsoidal level sets of $V$---does not.
Thus the statistical quantity $V$, the quadratic form of Chapter~II, and the
$(0,2)$ tensor of the present chapter are three descriptions of the same
invariant structure.

\medskip
\hrule
\bigskip

\bigskip
\hrule
\medskip

\textit{Note: OLAP Cubes and Tensor Contraction}

\medskip

In large-scale data analysis, observations are frequently organized along
multiple independent categorical dimensions such as time, location, system
component, or measurement type.
The resulting structure is commonly called an \emph{OLAP cube} (Online
Analytical Processing cube).

\medskip

Mathematically, such an object assigns numerical values to elements of a
Cartesian product
\[
D_1 \times D_2 \times \cdots \times D_n,
\]
where each set $D_i$ represents one observational dimension.
After enumerating these dimensions, the data can be written as a
multidimensional array
\[
F_{i_1 i_2 \dots i_n},
\]
which is precisely the coordinate representation of a tensor indexed by $n$
independent parameters.

\medskip

In practice, only a small subset of all possible index combinations is observed,
so the array is typically sparse.
This again reflects a distinction emphasized throughout this chapter: arrays
depend on chosen coordinates and storage conventions, while the underlying
multilinear structure is independent of representation.

\medskip

Analytical queries often aggregate data along selected dimensions.
For instance, combining observations while ignoring certain indices produces
quantities of the form
\[
G_{i_1 \dots i_k}
=
\sum_{j_1,\dots,j_m}
F_{i_1 \dots i_k j_1 \dots j_m}.
\]
From the tensorial viewpoint, this operation is a contraction: indices
corresponding to aggregated dimensions are summed over, reducing the order
of the tensor.

\medskip

Thus operations known operationally as filtering, grouping, or aggregation
correspond mathematically to projections of a high-dimensional tensor onto
lower-dimensional subtensors.
The OLAP cube therefore provides a concrete realization of tensorial structure
arising in modern data analysis.

\medskip

This example reinforces a recurring theme: tensor contraction appears wherever
multidimensional relations are summarized through systematic aggregation.

\medskip
\hrule
\bigskip